\chapter{Связь MLC и леммы о факторизации}
\noindent В главе устанавливается, что равномерная факторизация локальных компонент эквивалентна локальной связности множества Мандельброта. Компактностная редукция сводит конечное включение к отсутствию «плохих» точек в предельном множестве; обратное направление строит связную окрестность как пересечение вложенных компактных компонент. Тем самым условие равномерной факторизации является точным критерием MLC.

\section{Компактностная редукция факторизации}
\label{sec:factorization-compactness-reduction}

Для последующего сравнения предельного множества с конечными
приближениями зафиксируем $c\in\M$, $0<\delta<r$ и $N\in\N$ и
положим
\begin{equation}
 \label{eq:factorization-local-data}
 F_N(c,r):=O_N\cap\overline B(c,r),\qquad
 \mathcal B_N(c;r,\delta):=
 \bigl(F_N(c,r)\setminus C_N(c,r)\bigr)\cap\overline B(c,\delta).
\end{equation}
Множество $F_N(c,r)$ компактно и полуалгебраично. Поэтому оно имеет
конечное число компонент, каждая из которых компактна, и
$\mathcal B_N(c;r,\delta)$ компактно.

\begin{proposition}[Компактностное сведение факторизации]
\label{prop:factorization-compactness-reduction}
Для фиксированных $c\in\M$, $0<\delta<r$ и $N\in\N$ эквивалентны
условия
\begin{equation}
 \label{eq:factorization-finite-inclusion}
 \exists L\ge N,\qquad
 O_L\cap\overline B(c,\delta)\subseteq C_N(c,r)
\end{equation}
и
\begin{equation}
 \label{eq:no-bad-limit}
 \mathcal B_N(c;r,\delta)\cap\M=\varnothing.
\end{equation}
\end{proposition}
\begin{proof}
Из~\eqref{eq:factorization-finite-inclusion} и $\M\subseteq O_L$
следует~\eqref{eq:no-bad-limit}. Обратно, если
$\mathcal B_N(c;r,\delta)\cap\M=\varnothing$, вложенные компактные
множества $\mathcal B_N(c;r,\delta)\cap O_L$, $L\ge N$, не могут быть
все непустыми: их пересечение равно
$\mathcal B_N(c;r,\delta)\cap\M$. Для некоторого $L\ge N$ поэтому
$\mathcal B_N(c;r,\delta)\cap O_L=\varnothing$. Из $O_L\subseteq O_N$
и $\delta<r$ следует~\eqref{eq:factorization-finite-inclusion}.
\end{proof}

\section{Равномерный контроль компонент и локальная связность}

Обозначим через \(\operatorname{Disc}(S)\) дискретную категорию с
множеством объектов \(S\). Для фиксированных \(c\in\M\) и
\(0<\delta<r\) положим
\[
 \mathbf P_L(c,\delta):=
 \operatorname{Disc}\bigl(\pi_0(O_L\cap\overline B(c,\delta))\bigr),
 \qquad
 \mathbf C_N(c,r):=
 \operatorname{Disc}\bigl(\pi_0(O_N\cap\overline B(c,r))\bigr).
\]
При \(L\ge N\) включение
\[
 O_L\cap\overline B(c,\delta)
 \subseteq O_N\cap\overline B(c,r)
\]
переводит каждую компоненту связности левой части в единственную
компоненту правой части и задаёт функтор
\[
 F_{L,N}^{c;\delta,r}:
 \mathbf P_L(c,\delta)\longrightarrow\mathbf C_N(c,r).
\]
Включения при уточнении уровней также задают функторы
\[
 u_{L+1,L}:\mathbf P_{L+1}(c,\delta)\to\mathbf P_L(c,\delta),
\qquad
 v_{N+1,N}:\mathbf C_{N+1}(c,r)\to\mathbf C_N(c,r).
\]
Для \(L\ge N\) они удовлетворяют равенству
\[
 v_{N+1,N}\circ F_{L+1,N+1}^{c;\delta,r}
 =
 F_{L,N}^{c;\delta,r}\circ u_{L+1,L}.
\]
Равенство следует из того, что оба функториальных перехода посылают
компоненту более глубокого уровня в компоненту более грубого уровня,
содержащую её.

Определим функторы реализации
\[
 \mathcal Y_L:\mathbf P_L(c,\delta)\to\mathbf{CompHaus},
 \qquad
 \mathcal X_N:\mathbf C_N(c,r)\to\mathbf{CompHaus}
\]
отображением каждого объекта в соответствующую компактную компоненту;
компоненты компактных множеств замкнуты и потому компактны.
Включения компонент задают естественные преобразования
\[
 \mathcal Y_{L+1}\Rightarrow\mathcal Y_L\circ u_{L+1,L},\qquad
 \mathcal X_{N+1}\Rightarrow\mathcal X_N\circ v_{N+1,N},\qquad
 \mathcal Y_L\Rightarrow\mathcal X_N\circ F_{L,N}^{c;\delta,r}.
\]

Компонента \(C_N(c,r)\), содержащая \(c\), является выделенным объектом
\(\mathbf C_N(c,r)\). Условие \(\mathsf{UC}\) равносильно требованию,
чтобы для выбранных \(r,\delta\) при каждом \(N\) существовал \(L\ge N\),
для которого
\begin{equation}
 \forall A\in\pi_0(O_L\cap\overline B(c,\delta)),\qquad
 F_{L,N}^{c;\delta,r}(A)=C_N(c,r).
 \label{eq:component-functor-factorization}
\end{equation}
Иными словами, функтор от категории более глубоких локальных компонент
должен факторизоваться через полную подкатегорию, состоящую из одного
выделенного объекта. Коммутативность переходных функторов следует из
вложенности множеств; факторизация~\eqref{eq:component-functor-factorization}
есть дополнительное геометрическое утверждение.

\begin{theorem}[Эквивалентность равномерного контроля компонент и локальной связности]
\label{thm:uniform-component-control}
\[
 \mathsf{UC}\Longleftrightarrow\LC(\M).
\]
\end{theorem}
\begin{proof}
Сначала предположим \(\LC(\M)\). Зафиксируем \(c\in\M\) и
\(\varepsilon>0\). Локальная связность даёт связную окрестность \(V\) точки \(c\),
открытую в \(\M\), для которой
\[
 c\in V\subseteq \M\cap B(c,\varepsilon/2).
\]
Так как \(V\) открыто в \(\M\), существует \(d>0\) такое, что
\(\M\cap B(c,d)\subseteq V\). Положим
\[
 \delta:=\min\{d/2,\varepsilon/4\},\qquad r:=3\varepsilon/4,
 \qquad W:=\overline V^{\,\M}.
\]
Тогда \(0<\delta<r<\varepsilon\),
\(\M\cap\overline B(c,\delta)\subseteq V\), а \(W\) компактно и
связно, содержит \(c\) и лежит в
\(\M\cap\overline B(c,\varepsilon/2)\subseteq
\overline B(c,r)\). Поэтому \(W\subseteq C_N(c,r)\) для каждого
\(N\): это связное подмножество \(F_N(c,r)\), содержащее \(c\).
В частности,
\[
 \mathcal B_N(c;r,\delta)\cap\M=\varnothing
 \qquad(N\in\N).
\]
Предложение~\ref{prop:factorization-compactness-reduction} даёт для
каждого \(N\) индекс \(L\ge N\), для которого
\[
 O_L\cap\overline B(c,\delta)\subseteq C_N(c,r).
\]
Значит, выполнено \(\mathsf{UC}\).

Теперь предположим \(\mathsf{UC}\). Зафиксируем \(c\in\M\) и
относительную открытую окрестность \(U\) точки \(c\) в \(\M\).
Выберем \(\varepsilon>0\) так, чтобы
\(\M\cap B(c,\varepsilon)\subseteq U\), и применим
\eqref{eq:uniform-component-control}. Положим
\[
 C:=\bigcap_{N\in\N}C_N(c,r).
\]
Из \(O_{N+1}\subseteq O_N\) следует
\(C_{N+1}(c,r)\subseteq C_N(c,r)\). Множество
\(O_N\cap\overline B(c,r)\) компактно; его компонента \(C_N(c,r)\)
замкнута, поэтому компактна, связна и содержит \(c\). Теорема о
пересечении убывающей последовательности непустых компактных связных
множеств показывает, что \(C\) непусто, компактно и связно.
Кроме того, \(C\subseteq\bigcap_NO_N=\M\) и
\(C\subseteq\overline B(c,r)\).

Пусть \(y\in\M\cap\overline B(c,\delta)\). Для каждого \(N\) условие
\eqref{eq:uniform-component-control} даёт \(L\ge N\) с
\[
 O_L\cap\overline B(c,\delta)\subseteq C_N(c,r).
\]
Так как \(y\in\M\subseteq O_L\), имеем \(y\in C_N(c,r)\). Это верно
для каждого \(N\), следовательно,
\(\M\cap\overline B(c,\delta)\subseteq C\). Итак,
\[
 c\in\M\cap B(c,\delta)\subseteq C
 \subseteq\M\cap\overline B(c,r)
 \subseteq\M\cap B(c,\varepsilon)\subseteq U.
\]
Значит, \(C\) — связная окрестность \(c\), содержащаяся в \(U\).
Поскольку \(c\) и \(U\) произвольны, \(\M\) локально связно.
\end{proof}

Таким образом, реализация выделенных объектов вместе с функториальными
вложениями даёт связное множество
\(\bigcap_N C_N(c,r)\subseteq\M\). Факторизация
\eqref{eq:component-functor-factorization} обеспечивает, что это
множество содержит \(\M\cap\overline B(c,\delta)\), то есть является
окрестностью \(c\) в \(\M\). После этого предложение
\ref{prop:equicontinuous-selection} даёт последовательность отображений
с общим модулем непрерывности.

В частности, \(\mathsf{UC}\) влечёт \(\mathsf{EA}\) по
предложению~\ref{prop:equicontinuous-selection}, а значит, по
предложению~\ref{prop:diagonal-finite-scale} даёт требуемые конечные
оценки осцилляции при сколь угодно больших уровнях. Для конкретных
внешних множеств \(O_N\) завершение маршрута требует геометрической
факторизации~\eqref{eq:component-functor-factorization}; естественность
функторов задаёт согласованность переходов между уровнями.
